\pdfoutput=1
\documentclass[11pt]{article}
\usepackage[T1]{fontenc}
\usepackage{lmodern}
\usepackage{amsmath,amssymb,amsthm,mathtools,booktabs,array}
\usepackage[margin=1in]{geometry}
\usepackage[colorlinks=true,urlcolor=blue,citecolor=blue]{hyperref}
\hypersetup{pdftitle={Exact square and near-square Zarankiewicz numbers for K(3,3) through order 23},pdfauthor={Jason Davies}}
\newtheorem{theorem}{Theorem}
\newtheorem{lemma}[theorem]{Lemma}
\newtheorem{proposition}[theorem]{Proposition}
\newtheorem{corollary}[theorem]{Corollary}
\newcommand{\code}[1]{\nolinkurl{#1}}
\newcommand{\proofdir}[2]{\href{https://github.com/jasondavies/zarankiewicz-numbers/blob/5f5674e9d15a5abbc419501aabc8fda940c55d6d/#1}{#2}}
\emergencystretch=2em
\title{Exact square and near-square Zarankiewicz numbers\\for $K_{3,3}$ through order 23}
\author{Jason Davies\thanks{\href{mailto:jason@jasondavies.com}{\nolinkurl{jason@jasondavies.com}}}}
\date{November 27, 2026}
\begin{document}
\maketitle
\begin{abstract}
We determine the maximum number of ones in an $n\times n$ binary matrix
avoiding an all-ones $3\times3$ submatrix for $17\le n\le23$. The maxima
are $138,150,164,180,193,210,225$. We also determine the intervening
$n\times(n+1)$ maxima for $17\le n\le22$. The extremal $22\times22$
matrix is unique up to row and column permutations. For the upper bounds
we mark a line, enumerate the induced families of sets, and test which
families can carry the remaining entries. At $22\times23$, local bounds
on the degrees incident with a column give two incompatible counts of the
same incidences at weight 217. A matrix with 216 ones attains the bound.
At the next order we obtain $240\le z(24,24;3,3)\le243$.
\end{abstract}

\noindent\textbf{AI transparency statement.} The author directed an
AI-assisted investigation. AI systems developed proofs and verification
programs and drafted this article. The upper bounds are computer-assisted.
The scope of the checks is described in Section~\ref{sec:replay}.

\section{Introduction}
Let $z(m,n)$ be the largest number of ones in an $m\times n$ binary matrix
in which no three distinct rows have ones in the same three distinct columns.
Equivalently, it is the maximum number of edges in a $K_{3,3}$-free bipartite
graph with parts of sizes $m$ and $n$. Thus $z(m,n)=z(n,m)$. We use the
maximum-ones convention throughout. The least number of ones forcing a
forbidden submatrix is $z(m,n)+1$.

\begin{theorem}\label{thm:main}
The square and rectangular maxima are as follows.
\[
\begin{array}{c@{\qquad}rrrrrrr}
\toprule
n&17&18&19&20&21&22&23\\
\midrule
z(n,n)&138&150&164&180&193&210&225\\
z(n-1,n)&132&143&156&171&185&201&216\\
\bottomrule
\end{array}
\]
The extremal $22\times22$ matrix is unique up to independent
permutations of its rows and columns.
\end{theorem}

The first rectangular entry was determined by Wei~\cite{wei}. Four of the
seven square upper bounds follow immediately from the preceding rectangle
by averaging the weights after a row deletion.
One auxiliary value, $z(19,21)=177$, has two uses: an extra eight-one line
would be needed to exceed 185 on $20\times21$, and two such parallel lines
would be needed to exceed 193 on $21\times21$. In both cases removing the
lines would leave 178 on $19\times21$. In contrast, once the $22\times23$
bound is known, the upper bound of 225 on the next square follows by a
single averaging argument.

\paragraph{Earlier bounds.}
Table~4 of Collins, Riasanovsky, Wallace and Radziszowski~\cite{crwr}
gives upper bounds of 141, 148 and 156 on $17\times17$, $17\times18$ and
$18\times18$, respectively. Padhi~\cite{padhi} had already obtained the
lower bound of 138 on $17\times17$. Afrasyab~\cite{afra} gave
$132\le z(16,17)\le133$. Wei~\cite{wei} established $z(16,17)=132$ and
$z(16,18)=137$. Independent marked-line proofs for Wei's two values are
included in the accompanying repository~\cite{code}. The 137-one
witness there extends Hou's 132-one matrix on $15\times18$~\cite{hou}.
The provenance of the first entries is summarized below.
\begin{center}\small
\begin{tabular}{@{}lll@{}}
\toprule
Size & Earlier bound or exact value & Result used here\\
\midrule
$16\times17$ & $132$ exactly (Wei) & independent verification\\
$16\times18$ & $137$ exactly (Wei) & input for $17\times18$\\
$17\times17$ & $138$ (Padhi) to $141$ (Collins et al.) & upper bound $138$\\
$17\times18$ & upper bound $148$ (Collins et al.) & exact value $143$\\
$18\times18$ & upper bound $156$ (Collins et al.) & exact value $150$\\
\bottomrule
\end{tabular}
\end{center}

A consequence is the upper bound $b(2,2,3)\le17$: two $K_{2,2}$-free
colour classes on $17\times17$ have at most 74 cells each~\cite{crwr},
and $74+74+138<289$. Together with the known lower bound this recovers
the equality proved by Rowshan and Gholami~\cite{rg}.

\section{Marking a line}\label{sec:mark}
A \emph{line} is a row or a column. Its degree is its number of ones.
The weight of a matrix is the sum of its row degrees. Erasing a one
preserves avoidance, so an exclusion at one weight excludes every larger
weight on the same dimensions. If a matrix of weight $w$ has $m$ rows,
then each one survives $m-1$ of its $m$ single-row deletions. Consequently
\begin{equation}\label{eq:delete}
(m-1)w\le m\,z(m-1,n),\qquad d_r\ge w-z(m-1,n)
\end{equation}
for each row degree $d_r$. The transpose gives the column versions.

\begin{lemma}\label{lem:mark}
Mark $t$ columns in the support of one row. Restrict each other row to
these columns and regard the result as a block on $t$ points. Two blocks
intersect in at most two points. The load of point $j$, its number of
occurrences in the blocks, is $c_j-1$, where $c_j$ is its column degree.
If two rows share $u$ columns, the sum of degrees of those columns is
at most $2u+2(m-2)$ in a matrix with $m$ rows.
\end{lemma}
\begin{proof}
Three common points of two blocks give three common entries of their
rows and the marked row. For the last assertion the two rows contribute
$2u$ entries on the shared columns and each other row contributes at most
two.
\end{proof}

To enumerate the blocks, begin with all possible families of the largest
allowed size, often five. Adjoin smaller blocks in descending order.
A block of size three or more cannot repeat. A block of size at most two
can repeat, and the empty block accounts for a row missing the mark.
Permuting the points or the blocks gives an equivalent projection. We
canonicalize the bipartite incidence graph with points and blocks in
separate vertex classes using nauty~\cite{nauty}. Repeated blocks have
distinct vertices. In particular the enumeration must include the empty
family of five-blocks whenever the load bound permits it.

For a mark comprising the entire support of a row, the missing columns
avoid that row. On the remaining rows generate column supports $S$ of the
permitted sizes. For each pair of fixed columns $A,B$ require
$|S\cap A\cap B|\le2$. Upon choosing a column impose the same test on all
newly formed pairs. This checks exactly the triples involving a new
column. If the mark uses only part of the row, a missing column through
that row can use only blocks with pairwise intersection at most one.
The completion tests which exclude this possibility explicitly check how
many such blocks can be chosen.

\paragraph{Choosing a mark.}
Give column $j$ a weight $e_j$ and set $W_r=\sum_{j\sim r}e_j$, where
$j\sim r$ denotes a one. The identity
\begin{equation}\label{eq:weights}
\sum_r W_r=\sum_j c_j e_j
\end{equation}
can force a row with high incident weight even when ordinary averaging
only gives its degree. With $e_j=c_j-b$ and nonnegative weights, a
degree-$t$ row of incident weight $W_r$ gives a projection of total load
$(b-1)t+W_r$. Thus each extra unit forced by~\eqref{eq:weights} raises
the threshold in the block enumeration.

\section{The maxima through order 21}\label{sec:early}
Here and below the accompanying finite enumerations are part of the proof.
Appendix~\ref{app:files} identifies their inputs and entry points. We
describe the covering arguments since an empty list of completions gives
an upper bound only if every hypothetical matrix has a mark on the list.

\paragraph{The first square.}
At weight 139 on $17\times17$, deletion and the bounds of~\cite[Table~4]{crwr}
at heights 16, 15, 14 and 13 give minimum degree six and total deficit
below eight at most three on each side. Indeed there cannot be two sixes,
one six and two sevens, or four sevens on the same side. Six common
columns of two rows would need at least $48-3=45$ entries, exceeding
the capacity $12+2\cdot15=42$. Marking ten entries of a long line
now gives blocks of size at most five and total load at least 67.
For the five-blocks the $\{\pm1\}$ incidence vectors have squared norm
ten and mutual inner products at most $-2$. The squared norm of their
sum bounds their number by six. Enumerating the compatible fours and
the remaining blocks first excludes degrees above ten, then leaves
519 relabelling keys for a mark of degree ten. None extends.

For this argument the inputs at heights $16,15,14,13$ are, in order,
$133,126,118,110$ as upper bounds on seventeen columns. The six-row
case also uses the published uniqueness of the 128-one $16\times16$
matrix. A column of degree five in a 133-one core would have to extend
that square. Testing its $\binom{16}{5}$ possible supports in each
orientation excludes such an extension.

The remaining lines have degrees six through nine. If a six exists,
transpose as needed and delete a six-row. Otherwise delete a seven-row
if one exists. The respective $16\times17$ cores have 133 and 132 ones,
so each has a nine-row. Mark a nine-column if available and otherwise a
nine-row. With minimum degree eight on both sides of the original matrix,
mark a nine-column in the full square. These cases exhaust the possible
degrees. Completing the nine-point projections, including restoration
of the deleted row when present, excludes 139. The repository gives
two witnesses for 138, distinguished even up to transposition by their
maximum line degrees, ten and nine.

\paragraph{The next three rectangles and squares.}
At 144 on $17\times18$ a degree-nine column can be forced and no
column has greater degree. Project onto the nine rows containing the
mark or the eight rows missing it, according to the weight on each
side. Let $S$ denote the weight on the nine incident rows away from
the column. At $S\ge65$ the inside enumeration gives 3,671,498
projections. Otherwise the eight outside rows have at least 71 ones.
Their minimum degree is seven, eight or nine, giving respectively
7,298,758, 12,057,949 and 825 projections. These four cases give
23,029,030 in total. None
completes. Appending a five-one column to the second 138-one square gives
143. At 151 on $18\times18$ every degree is at least eight. Marking
eleven points on a largest line of degree at least eleven requires nine
five-blocks among seventeen blocks, with total load at least 77. No
such projection has the required loads. For maximum degrees ten and
nine there are respectively 978 and 1,821 projections of the required
types. In the latter case a nine-column meeting at least three nine-rows
is forced. None completes. The 150-one witness has three lines of
degree ten and fifteen of degree eight on each side.

At 157 on $18\times19$ the search marks a row after removing columns
of degree seven. The separate cases for zero and one such column and
the bounds on the smaller rectangles leave two through nine columns
to restore. A compatible single column need not extend to the required
family of columns. For example, in a cross-check with three columns
removed, 313 of 5,964 completed cores accepted one column and none
accepted three. The completed enumeration excludes 157 for each
number removed. Two witnesses give 156. One accepts an eight-one row
to give 164 on $19\times19$, attaining the bound
$\lfloor19\cdot156/18\rfloor$.

At 172 on $19\times20$, all degrees are at least eight, and an eight-row
and an eight-column cannot meet: deleting both would leave 157 ones on
$18\times19$. If there is an eight-row, its eight columns have degree
at least nine. It gives 114 possible marks, none of which completes.
Otherwise the rows have degrees $10,9,\ldots,9$. Marking ten rows
on any column of degree at least ten gives an empty catalogue. There
must then be a nine-column through the ten-row, and its 81 possible
projections all fail to complete. Thus $z(19,20)\le171$ and
$z(20,20)\le180$ by~\eqref{eq:delete}. For the lower bounds take all
translates of $\{0,2,3,6,7,8,9,11,15\}$ in $\mathbb Z/20\mathbb Z$
as rows, and delete a row for the rectangle. Direct triple intersection
gives the required avoidance.

\paragraph{One exhaustion for three values.}
At weight 178 on $19\times21$ let $s$ be the number of seven-columns.
Every column has degree at least seven, and deleting three seven-columns
would leave 157 on $19\times18$. Thus $0\le s\le2$. Remove all $s$.
The remaining column degrees have the form $8+e_j$ and the core row
degrees are at least eight, by the respective 18-row bounds 170, 163 and
156. The sum of the $e_j$ is $10+s$. For a row of degree $d$ the
projection has load $7d+W_r$. The two searches mark a row maximizing
$d+W_r$ or $2d+W_r$. Applying~\eqref{eq:weights} and the point-load
bounds gives the finite cases below. The counts include different
choices of mark for the same core.
\[
\begin{array}{lrrr}
\toprule
s&0&1&2\\
\midrule
\text{Marked cases}&155{,}847&174{,}743&184{,}800\\
\text{Extensions to weight 178}&0&0&0\\
\bottomrule
\end{array}
\]
The reductions at 186 on $20\times21$ and 194 on $21\times21$ force,
respectively, one eight-row and two parallel eight-lines (transposing if
necessary) whose removal gives exactly this 178-one case. This yields
the three upper bounds 177, 185 and 193. For the reverse inequalities,
start with the 193-one square and delete one or both of its eight-rows.
Details of the two score bounds and the reductions appear in the
\proofdir{z19x21-k33/README.md}{178-one exhaustion} and
\proofdir{z20x21-k33/weight-186.md}{eight-line reductions}.

\section{A 159-case proof at \texorpdfstring{$21\times22$}{21 by 22}}\label{sec:201}
\begin{proposition}
$z(21,22)=201$ and $z(22,22)=210$.
\end{proposition}
\begin{proof}
A column extends the 193-one square to 201. An additional row gives 210.
For the upper bounds,~\eqref{eq:delete} gives $z(21,22)\le202$ from
$z(21,21)=193$, and $z(20,22)\le193$ from $z(20,21)=185$.
Hence at weight 202 every row and column has degree at least nine.
Lemma~\ref{lem:mark} bounds pair intersections by five on both sides.

If a row has degree at least eleven, mark eleven of its entries. The
twenty blocks have total load at least 88 and size at most five. There
must be at least eight five-blocks. There are six compatible families
of at least eight five-blocks on eleven points up to relabelling.
For a family of size $m$, one needs at least $28-2m$ compatible
four-blocks, since the total load with $q$ fours is at most $60+2m+q$.
None of the six families admits them. Thus all rows have degree nine
or ten, with thirteen of degree ten.

Write $c_j=9+e_j$. Then $\sum e_j=4$. If one $e_j\ge2$, that column
must meet a ten-row, since it has degree at least eleven and there
are only eight nine-rows. Otherwise four columns have degree ten.
If every ten-row met at most one of the four, the eight nine-rows
would supply at least $40-13=27$ entries on them. By discrete convexity
the number of column triples on these rows would be at least
$5\binom33+3\binom43=17$, whereas the available capacity is
$2\binom43=8$. Either way a ten-row has $W_r\ge2$.

Mark that row. Its twenty blocks on ten points have point loads
from eight to twelve and total load at least 82. There are 159
canonical projections. The twelve missing columns avoid the marked
row. They have degrees from nine to eleven, because at most two excess
units remain outside the mark. Testing compatible supports as in
Section~\ref{sec:mark} gives zero completions to weight 202.
Finally $z(22,22)\le\lfloor22\cdot201/21\rfloor=210$.
\end{proof}

\paragraph{Classification at equality.}
In a 210-one square all degrees are at least nine. Give a column
weight $c_j-9$. A row of degree at least eleven has incident weight at
most one: otherwise its eleven heaviest points give one of 338
projections of load at least 90, none of which extends. The same
restriction holds on transposing. Thus a degree-$d$ row with $d\ge11$
requires at least $d-1$ nine-columns. Two such rows can share at
most five of these, and three can share at most two. Distributing
their required incidences among the nine-columns and applying the
same bounds on the other side leaves seven pairs of degree profiles
up to transposition. In each case~\eqref{eq:weights} forces a ten-row
of incident weight at least five or a nine-row of weight at least eight.

The ten-row gives 13,205 projections. Exhausting their completions
returns one square. To enumerate the nine-row marks, omit a point of
least degree, enumerate eight-point projections, then restore the point
on blocks meeting pairwise in at most one point. This gives 81 distinct
nine-point projections, all without completions. The resulting square
agrees up to relabelling with the 210-one witness. In either orientation
none of its $\binom{22}{7}$ proposed seven-point column extensions is
admissible. The classification and this last test are both used in
the next rectangle. See the \proofdir{z22-k33/unique.md}{classification
at weight 210} for the seven profiles and the reconstruction.

\section{The last rectangle}\label{sec:217}
Figure~\ref{fig:witness} gives a 225-one square. Its last row has nine
ones, so the other rows attain 216. We now exclude a $22\times23$ matrix
of weight 217. A row of degree nine will be called a nine-row, and
similarly for other degrees and for columns.

\begin{lemma}[Degree restrictions at 217]\label{lem:input}
Every row has degree at least nine. There are $k=1$ or $2$ eight-columns,
and all other columns have degree at least nine. Each eight-column
contains at most one nine-row. There are at least seven nine-rows.
A nine-row avoiding the eight-columns satisfies $\sum_{j\sim r}(c_j-9)\le6$.
A ten-column meets at least three nine-rows, and a column of degree
$c\ge11$ meets at least $c-5$.
\end{lemma}

We give the reductions behind this computational lemma, since using
only ordinary deletion bounds at this step would allow too many profiles.
A column of degree below seven would leave more than 210 on deletion.
Degree seven is impossible by the extension test on the unique square.
For the row bound one needs $z(21,23)\le208$. At weight 209 the
degree and projection tests force four or five eight-columns, say $q$,
and every nine-row avoids them. Write $b$ for the number of nine-rows.
On columns of degree at least ten give each entry weight $c-9$.
Those column weights sum to $q+2$, each column meets at least three
nine-rows, and each nine-row carries weight at most three. Hence
$b\ge q+2$. On each eight-column the sum of the incident row degrees
minus ten per incidence is at most three. A row of degree $10+s$,
$s>0$, needs at least $s+2$ eight-columns. Summing its excesses on
these columns gives $3(b-1)\le3q$: the positive row surplus above ten
sums to $b-1$, and no nine-row meets these columns. The two inequalities
are incompatible. The computations supplying these local bounds are
itemized in the \proofdir{z21x23-k33/no-209.md}{209-one exclusion}.
Deleting a row at weight 217 now proves the row bound.

On $22\times23$ an all-high-degree column side is ruled out by
\proofdir{z22x23-k33/min9.md}{weighted marks}. Deleting eight-columns
then leaves weights $209,201,193,185$ after one through four deletions.
The smaller rectangular bounds and the
\proofdir{z22x23-k33/eight-column-count.md}{eight-column count} exclude
five or more. The \proofdir{z22x23-k33/four-column-cover.md}{four-column
test} and the \proofdir{z22x23-k33/three-profile-proof.md}{three-column
test} exclude four and three. Deleting one eight-column leaves a
209-one square with at most one eight-row, since deleting an eight-row
of that square leaves 201 on $21\times22$ with minimum row degree nine.
This proves the asserted incidence bound on the eight-column.

Three further marked cases are needed. An eight-column with incident
row-degree sum at least 84 gives 900 projections. A ten-column with
two nine-rows and eight ten-rows gives 23,526. A nine-row avoiding
the eight-columns with incident excess exactly seven gives 15,766.
None completes. Together with the earlier high-degree marks these
tests give the excess bound of six, at least three nine-rows on a
ten-column, and at most five rows of degree above nine on any longer
column. If there are at most five nine-rows, all positive column excess
is on $10+k$ ten-columns. These require $3(10+k)$ incidences with the
nine-rows, which can supply at most thirty. For a nine-row touching
an eight-column, the latter capacity follows by deleting that column
and then the residual eight-row: the resulting 201-one rectangle
needs at least three eight-columns, forcing the original nine-row to
meet at least $4-k$ nine-columns.

At exactly six nine-rows the \proofdir{z22x23-k33/six-nine-rows.md}{incidence
count} forces all columns to have degree eight, nine or ten. There are
three units of positive row surplus above degree ten. If a ten-column
meets exactly three nine-rows and a longer row, the mark consists of
three nine-rows, six ten-rows and one eleven-row. A second longer row
would require one eleven-row and eight further rows of degree at least
ten outside this mark. None of its 129,226 projections admits even the
corresponding partial placement. For clarity, let $m'$ count nine-rows on a ten-column. With two
eight-columns, all these counts equal three. With one eight-column,
the sum of $m'-3$ over the eleven ten-columns is at most two.
A ten-column with $m'>3$ can hold at most $m'-2$ units of row excess
above ten. If at least two long rows exist, all their ten-column
incidences must now lie on these exceptional columns. They have room
for at most four such incidences, whereas two long rows together need
at least five: even counting multiplicities, they can use at most
seventeen positions outside the ten-columns. A single long row would
have degree thirteen and could meet only ten-columns with $m'\ge5$.
There are none if $k=2$. For $k=1$ there is at most one, forcing the
row to use all twelve columns outside the tens. A ten-row avoiding the
eight-column needs at least six nine-columns. Two rows could not share
those six, by Lemma~\ref{lem:mark} and their total degrees on those
columns. Details of these capacities and the six-row case are in the
\proofdir{z22x23-k33/seven-nine-rows.md}{six-nine-row exclusion}.
This completes the reductions used in Lemma~\ref{lem:input}.

\begin{lemma}[Column scores]\label{lem:score}
Give row $r$ weight $d_r-9$ and define $W_j=\sum_{r\sim j}(d_r-9)$.
At weight 217,
\[
\begin{array}{c|ccc}
c_j&8&9&\ge10\\ \hline
W_j\text{ at most}&11&9&7 .
\end{array}
\]
\end{lemma}
\begin{proof}
The 900-projection exclusion gives $W_j\le83-8\cdot9=11$ for an
eight-column. Any two columns share at most five rows. Indeed, six
shared rows need at least 54 ones. Their two columns and the other
21 columns supply at most $12+21\cdot2=54$. Equality forces every
one of the six rows to have degree nine and each remaining column to
meet exactly two of them. An eight-column would then meet either six
or two nine-rows, contradicting Lemma~\ref{lem:input}.

Each row has degree at most fifteen. If it meets $t$ columns of degree
at least nine, it shares at least $8t$ incidences on those columns
with the 21 other rows. No other row shares six such columns, by
Lemma~\ref{lem:mark} and $6\cdot9>12+20\cdot2$. Hence $8t\le105$ and
the at most two eight-columns give the bound.

For a nine-column of score at least ten mark all nine points. For a
column of degree at least ten and score at least eight mark its ten
heaviest incident rows. These retain at least eight units of score,
since weights are nonnegative integers. The 22 other columns give
blocks of size at most five. The loads lie between eight and fourteen.
Catalogue all block multisets with the indicated minimum total.
A ten-point mark of score at least nine with at least five
positive-weight points was already excluded by the earlier
\proofdir{z22x23-k33/min9.md}{ten-point test}.
\[
\begin{array}{rrrrr}
\toprule
\text{Points}&\text{Minimum load sum}&\text{Projections}&
  \text{Remaining}&\text{Passing}\\
\midrule
9&82&60{,}469&60{,}469&0\\
10&88&453{,}448&452{,}058&0\\
\bottomrule
\end{array}
\]

There are two tests. For a column of degree greater than ten, a further
row on the column needs eight of the outside columns. Their blocks must
meet pairwise in at most one marked point. No remaining ten-point mark
admits an eight-clique in this compatibility graph. For a full mark of
$p=9$ or $10$ points, write $M$ for the total of the marked row degrees,
and let $h$ of those rows have degree nine. The $22-p$ outside rows
need a gain of $217-M-9(22-p)$ above their baseline of nine each.
At most $22-p-\max(0,7-h)$ of them can supply gain. Enumerate the
supports of degrees ten to fifteen on the 22 outside columns. Choose
compatible longer rows in decreasing degree order, then ten-rows,
testing each new triple. Two equal supports of size ten would share at
least eight columns of degree at least nine, already impossible.
A completion must attain the required gain within the available number
of rows. No projection in either row of the table passes this test.
\end{proof}

\begin{proposition}\label{prop:last}
$z(22,23)=216$ and $z(23,23)=225$.
\end{proposition}
\begin{proof}
Suppose there are 217 ones. Write $a$ for the number of nine-rows and
$m$ for the number of columns of degree at least ten. On the latter
columns put $e_j=c_j-9$, $E=\sum e_j=10+k$ and $D=\sum(e_j-1)=E-m$.
The rows above degree nine have degrees $10+s_r$, where $s_r\ge0$ and
$\sum s_r=a-3$. Summing the scores of Lemma~\ref{lem:score} first by
rows and then by columns gives
\begin{align*}
\sum_j W_j
 &=10(22-a)+11\sum s_r+\sum s_r^2\ \ge\ 184+2a,\\
\sum_j W_j
 &\le11k+9(23-k-m)+7m=207+2k-2m .
\end{align*}
Here $s_r^2\ge s_r$. Integrality gives $a+m\le11+k$, or $a\le D+1$.

Now let $T$ count incidences between the nine-rows and the $m$ columns,
each incidence counted with weight $e_j$. A column of weight one
contributes at least three. A column of weight $e\ge2$ meets at least
$e+4$ nine-rows, and
\[
 e(e+4)-3e-6(e-1)=(e-2)(e-3)\ge0.
\]
Hence $T\ge3E+6D$. A nine-row avoiding the eight-columns contributes
at most six. At most $k$ nine-rows meet an eight-column, and each can
contribute at most $E\le12$. Thus $T\le6a+6k\le6D+6+6k$.
Together these require $3(10+k)\le6+6k$, so $k\ge8$. This
contradicts $k\in\{1,2\}$. The witness attains 216. Applying
\eqref{eq:delete} gives $z(23,23)\le\lfloor23\cdot216/22\rfloor=225$,
also attained.
\end{proof}

% The following figure is drawn directly from the witness bits. Rows run
% top to bottom and columns left to right. No external graphics are needed.
\makeatletter
\newcount\gridx
\newcommand{\gridrow}[2]{\gridx=0\relax
  \@tfor\digit:=#2\do{\ifnum\digit=1\relax
    \put(\the\gridx,#1){\rule{\unitlength}{\unitlength}}\fi
    \advance\gridx by 1\relax}}
\makeatother
\begin{figure}[htbp]
\centering
{\setlength{\unitlength}{4pt}
\begin{picture}(21,21)
\gridrow{20}{110101000001000101010}
\gridrow{19}{100100001111110001100}
\gridrow{18}{001110000101101110010}
\gridrow{17}{011010010001100001101}
\gridrow{16}{000011000011011000111}
\gridrow{15}{110011001100101100001}
\gridrow{14}{001001111001011101000}
\gridrow{13}{100000110011100110001}
\gridrow{12}{111110010010011000000}
\gridrow{11}{111000100101010000011}
\gridrow{10}{010100100000011111101}
\gridrow{9}{010000011101001010100}
\gridrow{8}{010001010110110011010}
\gridrow{7}{101000001010001011011}
\gridrow{6}{101011100100000011100}
\gridrow{5}{010111101011000010000}
\gridrow{4}{011000101010100100110}
\gridrow{3}{100010011000010110110}
\gridrow{2}{000110111100000001011}
\gridrow{1}{100101110000101000110}
\gridrow{0}{001101010110000100101}
\end{picture}}\qquad {\setlength{\unitlength}{4pt}
\begin{picture}(23,23)
\gridrow{22}{01000111010010001001011}
\gridrow{21}{11001100110010110000100}
\gridrow{20}{01101001001110000010110}
\gridrow{19}{11111011000000100100001}
\gridrow{18}{00111000010110111001000}
\gridrow{17}{01010000001000111010111}
\gridrow{16}{01100000100011010111001}
\gridrow{15}{00001110000100100111101}
\gridrow{14}{00011100101010001100010}
\gridrow{13}{00100101101100110000011}
\gridrow{12}{01000001110101101110000}
\gridrow{11}{10010010110110000010011}
\gridrow{10}{01011110000101010000010}
\gridrow{9}{00011001111001000001101}
\gridrow{8}{10100000010001100101110}
\gridrow{7}{10000001000111011100101}
\gridrow{6}{10001011100000011011010}
\gridrow{5}{10101100011001001010001}
\gridrow{4}{10010101001011100011000}
\gridrow{3}{00110111010000010110100}
\gridrow{2}{11000010011100010101000}
\gridrow{1}{11110100100100001001100}
\gridrow{0}{00100010101011101000100}
\end{picture}}
\caption{Extremal squares of weights 193 (left) and 225 (right).
Removing the two eight-rows of the first square gives 177 on $19\times21$.
Removing the last row of the second gives 216 on $22\times23$.
Each black cell is a one.}\label{fig:witness}
\end{figure}

\section{The next order}\label{sec:24}
There is a short cyclic construction one order further. In
$\mathbb Z/24\mathbb Z$ take the translates of
$\{0,1,2,3,5,8,11,13,17,18\}$. Every triple of translates has
intersection of size at most two. This gives 240 ones on the square
and 230 after deleting any row. Column deletion gives
$z(23,24)\le\lfloor24\cdot225/23\rfloor=234$.

\begin{proposition}
$z(24,24)\le\max\{240,z(23,24)+9\}\le243$.
\end{proposition}
\begin{proof}
Above weight 240 some row has at least eleven entries. Suppose every
column has degree at least ten and mark eleven entries of that row.
A six-point block would require 60 incidences on its six columns,
above their capacity $12+22\cdot2=56$. Thus the 23 blocks have sizes
at most five, and total load at least $11\cdot9=99$. Write $m$ and $q$
for the number of fives and fours. Counting first relative to four,
then to three, gives $m\ge7$ and $q\ge30-2m$. There are 21 compatible
families with at least seven fives on eleven points up to relabelling.
For each, enumerate the choices of $30-2m$ pairwise compatible fours.
Only $m=7$ contributes, so all 23 blocks are specified. Of the
18 choices on the representative families, thirteen have minimum point
load eight and five have minimum seven. None reaches nine. Hence a
column of degree at most nine exists. Delete it and use transposition.
\end{proof}

At weight 234 on $23\times24$ every column has degree at least nine
and at least six have degree exactly nine. Each such deletion would
leave an extremal 225-one square. In both orientations of the stored
square exactly one five-point column extension is admissible, and
none of size six is. A 234-one rectangle would therefore require a
different 225-one square. Whether another 225-one square admits such an extension remains open in this analysis.

\section{Computational checks}\label{sec:replay}
For a lower bound, encode every row as an integer with a bit for each
column. An independent witness check consists of summing the bit counts
of the rows, then checking
$\operatorname{popcount}(r_i\mathbin{\&} r_j\mathbin{\&}r_k)\le2$ for
each $i<j<k$. For the preparation of this article all 18 stored witness
files were checked in both orientations from a clean export of the
revision in~\cite{code}. The eight-orbit construction of the 225-one
square can also be regenerated with \code{z23-k33/construct.py}. Clean-export replays for this article
also recovered the 171-one and 201-one bounds, the unique 210-one square,
and the two final catalogues excluding weight 217. Other enumeration
counts are reported from the individual proof records in~\cite{code}.

For an upper bound there are three separate claims to check: a line or
deletion of the asserted kind must exist, the catalogue must contain
every possible projection at that mark, and the completion test must
retain any valid completion. The numerical counts in the tables refer
to the indicated stage of the search. A positive control for the last
column-score calculation hides rows in the 216-one witness and recovers
a compatible placement. The final two exclusions can be run with
\begin{verbatim}
cd z22x23-k33
python3 verify_217.py --jobs 4
\end{verbatim}
This command checks the two rows of the table in Lemma~\ref{lem:score}.
The directory's proof index records the separate verifications on
which the degree restrictions depend. The final proof of
$z(21,23)\le208$ uses incidence counts, so the large DRAT archives for
the earlier individual four-column case are unnecessary for
Theorem~\ref{thm:main}. Run the Python drivers with assertions enabled
(omit \code{-O} and \code{-OO} and leave \code{PYTHONOPTIMIZE} unset).
The native replays use Python 3.10 or later, a C++20 compiler and 64-bit
nauty. The degree-profile replay for the three-eight-column case also uses
OR-Tools, SciPy and NumPy. It certifies 298 exclusions among 302 profile
pairs. These are program-checked finite searches. Their
correctness depends on the coverage arguments, the enumerators and the
completion routines.

\appendix
\section{Guide to the finite computations}\label{app:files}
All paths below are relative to the repository in~\cite{code}. The
linked notes specify the degree inequalities, branching rules,
multiplicities, and counts for their respective computations. Parameters
and counts refer to the same revision.

\begin{center}\small
\begin{tabular}{@{}lll@{}}
\toprule
Target or input & Computational step & Replay or proof index\\
\midrule
$17\times17$ & Six-, seven-, and nine-line cases
 & \proofdir{z17-k33/verify_exact.py}{\code{z17-k33/verify_exact.py}}\\
$17\times18$ & Both sides of a nine-column
 & \proofdir{z17x18-k33/verify_exact.py}{\code{z17x18-k33/verify_exact.py}}\\
$18\times18$ & Marks of size 9, 10, 11
 & \proofdir{z18-k33/verify_exact.py}{\code{z18-k33/verify_exact.py}}\\
$18\times19$ & Joint seven-column extension
 & \proofdir{z18x19-k33/verify_exact.py}{\code{z18x19-k33/verify_exact.py}}\\
$19\times20$ & Marks in both orientations
 & \proofdir{z19x20-k33/verify_exact.py}{\code{z19x20-k33/verify_exact.py}}\\
$19\times21$ & Three counts of seven-columns
 & \proofdir{z19x21-k33/verify_weight_178.py}{\code{z19x21-k33/verify_weight_178.py}}\\
$20\times21$, $21\times21$ & Reductions to 178
 & \proofdir{z20x21-k33/weight-186.md}{\code{z20x21-k33/weight-186.md}}\\
$21\times22$ & Marked ten-row
 & \proofdir{z21x22-k33/verify_exact.py}{\code{z21x22-k33/verify_exact.py}}\\
$22\times22$ & Census at 210
 & \proofdir{z22-k33/verify_unique.py}{\code{z22-k33/verify_unique.py}}\\
$21\times23$ & Exclusion of 209
 & \proofdir{z21x23-k33/no-209.md}{\code{z21x23-k33/no-209.md}}\\
$22\times23$ & Number of eight-columns
 & \proofdir{z22x23-k33/README.md}{\code{z22x23-k33/README.md}}\\
 & Local degree restrictions
 & \proofdir{z22x23-k33/verify_217_loads.py}{\code{z22x23-k33/verify_217_loads.py}}\\
 & Exclusion of six nine-rows
 & \proofdir{z22x23-k33/verify_six_nines.py}{\code{z22x23-k33/verify_six_nines.py}}\\
 & Final column-score bounds
 & \proofdir{z22x23-k33/verify_217.py}{\code{z22x23-k33/verify_217.py}}\\
\bottomrule
\end{tabular}
\end{center}

For the small calculation in Section~\ref{sec:24} the following
pseudocode uses \code{z17x18-k33/fives.cpp} and the clique iterator in
\code{z17x18-k33/verify.py}, both in the cited revision. Compatibility
means intersection of size at most two. The outer loop visits each
family size, including the nonmaximal families.
\begin{verbatim}
for F in canonical compatible families of at least seven
         5-subsets of an 11-set:
    P = all 4-subsets compatible with every member of F
    for Q in cliques(P, size = 30 - 2*len(F)):
        assert len(F) == 7 and len(Q) == 16
        assert some point lies in fewer than 9 members of F union Q
\end{verbatim}

\section*{Declaration of AI-generated research and authorship}
Jason Davies initiated the investigation, selected the sequence of
square and near-square cases, and directed the computational work.
AI systems working under his direction developed mathematical reductions,
wrote enumeration and checking programs, and produced this manuscript
draft. The programs and the scope of their checks are described above
so that the calculations can be examined independently.

\begin{thebibliography}{9}
\bibitem{crwr}
Alex F. Collins, Alexander W. N. Riasanovsky, John C. Wallace and
Stanis\l aw P. Radziszowski.
Zarankiewicz numbers and bipartite Ramsey numbers.
\emph{Journal of Algorithms and Computation} 47 (2016), 63--78.
\href{https://www.cs.rit.edu/~spr/PUBL/crwr16.pdf}{Author's copy}.

\bibitem{padhi}
Jaideep Padhi.
New exact values and improved lower bounds for Zarankiewicz numbers
$z(m,n;3)$. Preprint, 2026.
\href{https://papers.ssrn.com/sol3/papers.cfm?abstract_id=6960039}{SSRN 6960039}.

\bibitem{afra}
Koyar Afrasyab. Exact Zarankiewicz values on two finite frontier slices.
Preprint, 2026. \href{https://arxiv.org/abs/2608.08154}{arXiv:2608.08154}.

\bibitem{hou}
Shengteng Hou. Seven exact finite Zarankiewicz numbers from a single
$13\times18$ core. Preprint, 2026.
\href{https://arxiv.org/abs/2608.08549}{arXiv:2608.08549}.

\bibitem{wei}
Jiuzeng Wei. The Zarankiewicz numbers $z(16,17;3,3)=132$ and $z(16,18;3,3)=137$,
and the bound $z(16,19;3,3)\le143$. Research data, 2026.
\href{https://doi.org/10.5281/zenodo.24553293}{Zenodo 24553293}.

\bibitem{rg}
Yaser Rowshan and Mostafa Gholami. A proof of a conjecture on Ramsey
numbers $B(2,2,3)$. 2021.
\href{https://arxiv.org/abs/2108.03572}{arXiv:2108.03572}.

\bibitem{nauty}
Brendan D. McKay and Adolfo Piperno.
Practical graph isomorphism, II.
\emph{Journal of Symbolic Computation} 60 (2014), 94--112.
\href{https://doi.org/10.1016/j.jsc.2013.09.003}{doi:10.1016/j.jsc.2013.09.003}.

\bibitem{code}
Jason Davies. Zarankiewicz numbers. Explicit matrices, proof notes and
verification sources, 2026. Revision \code{5f5674e}.
\url{https://github.com/jasondavies/zarankiewicz-numbers/tree/5f5674e9d15a5abbc419501aabc8fda940c55d6d}.
\end{thebibliography}
\end{document}
